%% sections/analysis.tex
\section{Analysis and Discussion}
\label{sec:analysis}

% ---- Quality vs Speed Trade-off ----
\subsection{Quality--Speed Trade-off}
\label{sec:analysis:quality_speed}

Our benchmark reveals a fundamental trade-off between rendering quality
and speed. Methods that improve anti-aliasing (e.g., Mip-Splatting
\cite{yu2024mipsplatting}) achieve higher PSNR but at the cost of
rendering FPS. Conversely, compression-focused methods (e.g.,
Scaffold-GS \cite{yu2024scaffold}) reduce memory while maintaining
competitive speed.

\todo{Add scatter plot: PSNR vs FPS, with method labels.}

% ---- Memory Analysis ----
\subsection{Memory Efficiency}
\label{sec:analysis:memory}

The memory footprint of 3DGS varies significantly across methods.
The original 3DGS \cite{kerbl3Dgaussians} requires substantial storage
due to unpruned Gaussian growth. Scaffold-GS \cite{yu2024scaffold}
achieves \todo{X\%} reduction through anchored representation.

% ---- Geometric Accuracy ----
\subsection{Geometric Accuracy}
\label{sec:analysis:geometry}

2DGS \cite{huang20242dgs} demonstrates superior geometric accuracy
on the Drum Tower Dataset, benefiting from its surface-aligned
representation. This is particularly evident in scenes with thin
structures and complex occlusions.

\todo{Add normal map comparison and depth error visualization.}

% ---- Failure Cases ----
\subsection{Failure Cases and Limitations}
\label{sec:analysis:failure}

Despite impressive progress, 3DGS methods still struggle with:
\begin{itemize}[leftmargin=*]
    \item \textbf{Texture-less regions}: Large flat surfaces lack features
          for reliable optimization.
    \item \textbf{Specular materials}: View-dependent appearance is
          challenging for discrete Gaussian primitives.
    \item \textbf{Extreme lighting}: High dynamic range scenes exceed
          the representational capacity of spherical harmonics.
    \item \textbf{Thin structures}: Wire-like geometry is often
          under-represented.
\end{itemize}

% ---- Future Directions ----
\subsection{Future Directions}
\label{sec:analysis:future}

Based on our analysis, we identify the following promising directions:
\begin{enumerate}[leftmargin=*]
    \item \textbf{Scene-level understanding}: Integrating semantic
          features into 3DGS for downstream tasks.
    \item \textbf{Physics-aware representations}: Incorporating physical
          priors for more realistic rendering.
    \item \textbf{Edge deployment}: Optimizing 3DGS for mobile and
          embedded devices.
    \item \textbf{Unified frameworks}: Combining reconstruction,
          editing, and generation in a single 3DGS framework.
\end{enumerate}
